% from main.tex
%%%%%%%%%%  scan idx 518  |  folio 505  %%%%%%%%%%
présentation finie).

Ces propriétés étant locales sur $X$ et $S$, on peut les supposer affines; soient $S=\Spec(A)$, $X=\Spec(B)$. On a alors:
\[
\operatorname{tor.dim}_B(A)=\operatorname{dim.proj}_B(A)=\sup_{s\in S}\operatorname{dim.proj}_{B\otimes_A k(s)}(k(s))
\]
d'après EGA IV 12.3.2 où on peut remplacer l'hypothèse de présentation finie sur $B$ par des hypothèses noethériennes. Or $k(s)$ est isomorphe au corps résiduel de $\mathcal O_{X_s,g(s)}$, et d'après EGA $0_{\mathrm{IV}}$ 17.2.5 on a:
\[
\operatorname{dim.proj}_{B\otimes_A k(s)}(k(s))=
\operatorname{dim.proj}_{\mathcal O_{X_s,g(s)}}(k(s)).
\]
Il résulte alors de EGA $0_{\mathrm{IV}}$ 17.2.11 que cette dernière est finie si et seulement si $\mathcal O_{X_s,g(s)}$ est de dimension cohomologique finie, c'est-à-dire régulier.

\textbf{Proposition 1.6.} Soient $S$ un schéma noethérien, $E$ un $\mathcal O_S$-Module localement libre de type fini, $X$ le fibré vectoriel associé à $E$, et $f$ le morphisme structural de $X$. Alors \(f^*:K_\bullet(S)\longrightarrow K_\bullet(X)\) est un isomorphisme.

Comme $f$ est lisse, et admet une section, à savoir la section nulle, il suffit pour prouver 1.6 de montrer que pour tout $s\in S$, $K_\bullet(X_s)\simeq \mathbb Z$ (1.4). On est ainsi ramené au cas où le schéma de base est un corps.

Dans ce cas, $E$ est somme directe de sous-modules de rang $1$, ce qui permet de décomposer $f$ en une suite de fibrations de rang $1$, et donc de se ramener au cas où $X$ est de rang $1$, soit $X=\Spec(k[T])$. La proposition résulte alors du lemme suivant:

